\begin{abstract}
Building on recent advances such as BitNet~\cite{bitnet}, the landscape of Large Language Models (LLMs) is shifting toward an era dominated by 1-bit LLMs. In this paper, we propose a novel 1-bit LLM architecture, denoted as \textbf{\text{BitNet b1.58}}, where each individual parameter (or weight) is drawn from the ternary set \{-1, 0, 1\}. This architecture is able to achieve parity with its full-precision (FP16 or BF16) Transformer counterparts of the same model scale and trained on an equivalent number of tokens, both in terms of perplexity and downstream task results, while offering significant improvements in cost efficiency regarding latency, memory usage, throughput, and energy expenditure. Notably, the 1.58-bit LLM introduces a distinct scaling law and a new methodology for training high-performing yet resource-efficient LLMs. In addition, it ushers in a novel computational paradigm, presenting opportunities for the development of bespoke hardware tailored to optimize 1-bit LLMs.
\end{abstract}

\section{The Era of 1-bit LLMs}

Artificial intelligence research has witnessed the swift expansion of both the size and the capability of Large Language Models (LLMs) in recent years. These advancements have enabled exceptional outcomes on various natural language processing benchmarks. However, the escalating scale of these models introduces significant barriers to practical deployment, exacerbated by concerns over their environmental footprint and economic cost stemming from increased energy requirements.

To mitigate these issues, one popular strategy is post-training quantization, which yields low-bit representations for efficient inference~\cite{smoothquant, gptq, quip, quip_sharp}. By decreasing the numerical precision of model weights and activations, this method drastically cuts down the memory footprint and compute demand of LLMs. Industry trends have progressed from 16-bit formats toward even more compact representations, such as 4-bit implementations~\cite{gptq, awq}. Despite its prevalence in industry applications, post-training quantization does not fully address performance trade-offs and remains inherently suboptimal.

Recent advances in 1-bit neural network architectures, such as BitNet~\cite{bitnet}, demonstrate a compelling approach to lowering the computational demands of large language models (LLMs) without sacrificing accuracy. Conventional LLMs typically utilize 16-bit floating-point formats (e.g., FP16 or BF16), and their main operations revolve around matrix multiplications, where floating-point addition and multiplication account for the predominant computational expense. BitNet, on the other hand, performs matrix multiplications using only integer addition, resulting in substantial energy savings for LLM computations. Since energy consumption constitutes a critical limitation in the computational throughput of many hardware accelerators, this improvement translates not only to reduced power draw but also to potential speedup in processing.

Apart from the computation itself, inference also incurs costs when moving model parameters from off-chip DRAM to the accelerator's on-chip memory (such as SRAM). Although efforts have been made to increase SRAM size to enhance throughput, such measures are significantly costlier compared to DRAM. 1-bit LLMs offer a stark reduction in memory requirements regarding both storage size and bandwidth, especially relative to models operating in full precision. This reduction minimizes both the latency and the cost of retrieving parameters from DRAM, leading to faster, more resource-efficient inference.

In this paper, we present a new 1-bit LLM variant, designated as \textbf{\text{BitNet b1.58}}, in which every model parameter is selected from the ternary set \{-1, 0, 1\}. By incorporating an additional value (0) into the original 1-bit BitNet, \text{BitNet b1.58} achieves a representation of 1.58 bits in binary terms. All principal advantages of the original 1-bit BitNet are preserved in \text{BitNet b1.58}, such as a computation strategy that largely eliminates the need for multiplication in matrix operations and is amenable to aggressive optimization. Moreover, \text{BitNet b1.58} maintains comparable energy efficiency and offers improvements over FP16 LLMs in memory usage, data throughput, and inference latency. Additionally, this variant introduces two further key benefits. First, its expressivity is enhanced owing to explicit feature filtering enabled by the inclusion of zero-valued weights, significantly boosting 1-bit LLM performance. Second, experimental results confirm that \text{BitNet b1.58} can attain parity with full precision (FP16) baselines regarding perplexity and downstream task performance, at least at the 3B parameter scale and under matching experimental setups (including model configuration, token count, etc.).

\section{BitNet b1.58}

\text{BitNet b1.58} builds upon the foundational BitNet model, a Transformer variant in which the standard \emph{nn.Linear} layers are substituted with \emph{BitLinear} layers. This model is initialized and trained from the beginning, employing 1.58-bit precision for weights and 8-bit for activations. Relative to its predecessor, BitNet b1.58 incorporates several adjustments, which we outline below.

\paragraph{Quantization Function.} We utilize an \emph{absmean} quantization approach to restrict weight values to -1, 0, or +1. The procedure involves dividing the weight matrix by its mean absolute value before mapping each entry to the closest integer among \{-1, 0, +1\}:

\begin{equation}
    \widetilde{W} = \mathrm{RoundClip}(\frac{W}{\gamma + \epsilon}, -1, 1),
\end{equation}

\begin{equation}
    \mathrm{RoundClip}(x, a, b) = \max(a, \min(b, \mathrm{round}(x))),
\end{equation}

\begin{equation}
    \gamma = \frac{1}{nm}\sum_{ij} |W_{ij}|.
\end{equation}

The same strategy is applied to activation quantization as in the original BitNet, except we omit the step of rescaling activations to $[0, Q_b]$ prior to nonlinear transformation. Instead, all activations are normalized to the interval $[-Q_b, Q_b]$ on a per-token basis, eliminating zero-point quantization. This adjustment streamlines both code implementation and system-level deployment, and our experiments reveal it introduces negligible impact on overall performance.

\paragraph{LLaMA-alike Components.}

LLaMA~\cite{llama, llama2} has become the standard reference architecture for many open-source large language models. To better align with community standards, the architecture of \text{BitNet b1.58} is modeled after LLaMA's design choices. This includes adopting RMSNorm~\cite{rmsnorm}, SwiGLU~\cite{swiglu}, and rotary embeddings~\cite{rope}, as well as eliminating bias terms throughout. Such design decisions enable seamless integration of \text{BitNet b1.58} into widely used open-source frameworks (such as Huggingface, vLLM~\cite{vllm}, and llama.cpp\footnote{\url{https://github.com/ggerganov/llama.cpp}}) with minimal modification required.

\section{Results}

We conducted a comparison between \text{BitNet b1.58} and our FP16 reproduction of \text{LLaMA LLM} across several model scales. To standardize the experimental conditions, both sets of models were pre-trained on the RedPajama dataset~\cite{redpajama} for a total of 100 billion tokens. Their zero-shot capabilities were then assessed on diverse language understanding benchmarks, namely ARC-Easy~\cite{arc}, ARC-Challenge~\cite{arc}, Hellaswag~\cite{hellaswag}, Winogrande~\cite{winoGrande}, PIQA~\cite{piqa}, OpenbookQA~\cite{openbookqa}, and BoolQ~\cite{boolq}. Additionally, validation perplexities were computed on WikiText2~\cite{wikitext2} and C4~\cite{c4} to further measure model performance.

For resource efficiency metrics, we evaluated both runtime GPU memory utilization and inference latency for \text{LLaMA LLM} and \text{BitNet b1.58}. Our measurements were obtained using the FasterTransformer\footnote{\url{https://github.com/NVIDIA/FasterTransformer}} library, which provides efficient GPU-optimized inference routines for large language models. For \text{BitNet b1.58}, the Ladder~\cite{ladder} 2-bit kernel was incorporated. We reported the time required to generate each output token, reflecting the primary cost driver during inference.

As shown in Table~\ref{tab:ppl}, a summary of perplexity and resource metrics for both \text{BitNet b1.58} and \text{LLaMA LLM} is presented. Notably, \text{BitNet b1.58} approaches the perplexity of full-precision \text{LLaMA LLM} at the 3B parameter scale, with a 2.71$\times$ increase in inference speed and a 3.55$\times$ reduction in GPU memory use. Furthermore, \text{BitNet b1.58} at 3.9B parameters surpasses \text{LLaMA LLM} 3B both in speed (2.4$\times$ faster), memory consumption (3.32$\times$ lower), and delivers markedly improved performance.

Table~\ref{tab:zero-shot} provides a detailed breakdown of zero-shot accuracy scores across evaluation tasks. The assessments were conducted using the \emph{lm-evaluation-harness}\footnote{\url{https://github.com/EleutherAI/lm-evaluation-harness}} framework. Results indicate that the performance differential between \text{BitNet b1.58} and \text{LLaMA LLM} diminishes as model size increases. Crucially, \text{BitNet b1.58} is able to equal the accuracy of its full-precision counterpart at the 3B scale and above. Consistent with perplexity outcomes, final task performance demonstrates that \text{BitNet b1.58} 3.9B exceeds the results of \text{LLaMA LLM} 3B, while demanding fewer computational resources in terms of memory and latency. This indicates that \text{BitNet b1.58} offers a Pareto-efficient improvement over state-of-the-art LLM architectures.

\paragraph{Memory and Latency}

We extended our experiments by increasing the model sizes to 7B, 13B, and 70B, and analyzed the corresponding computational costs. The results presented in Figure~\ref{fig:latency-memory-energy} indicate that both latency and memory usage benefit from scaling, with acceleration improving as model size increases. Notably, \text{BitNet b1.58} at 70B exhibits a 4.1$\times$ reduction in latency compared to the \text{LLaMA LLM} reference model. This advantage arises from the scaling of time required by \emph{nn.Linear}, which increases with larger architectures. A parallel pattern is observed for memory consumption: since the embeddings are retained in full precision, their share of total memory decreases as the models grow. Both latency and memory metrics were gathered using a 2-bit kernel, suggesting that there is further potential for optimization and cost reduction.

\paragraph{Energy}

In addition, we calculated the energy demands associated with arithmetic operations for both \text{BitNet b1.58} and \text{LLaMA LLM}, with an emphasis on matrix multiplication due to its dominant role in LLM resource usage. Figure~\ref{fig:energy} breaks down the components contributing to total energy expenditure. For \text{BitNet b1.58}, the bulk of energy use is from INT8 addition operations, whereas \text{LLaMA LLM} is characterized by both FP16 addition and FP16 multiplication. Drawing upon the energy models described in~\cite{energycost, pokebnn}, it is demonstrated that \text{BitNet b1.58} achieves a 71.4-fold reduction in matrix multiplication energy consumption on 7nm hardware. We additionally measured the total end-to-end energy usage for scenarios with 512 tokens. Our findings indicate that the energy efficiency of \text{BitNet b1.58} relative to the FP16 \text{LLaMA LLM} baseline improves further as the model size increases. This improvement is attributed to the higher proportion of cost incurred by \emph{nn.Linear} layers in larger models, while the resource demands from other network components diminish in significance.

\paragraph{Throughput}

We assess the throughput of \text{BitNet b1.58} in comparison to the \text{LLaMA LLM} with 70B parameters, utilizing two 80GB A100 GPUs and employing pipeline parallelism~\cite{gpipe} to enable the deployment of the \text{LLaMA LLM} 70B on the available hardware. The batch size was increased incrementally until the GPUs reached their memory capacity, with each input sequence being of length 512. According to Table~\ref{tab:throughput}, the \text{BitNet b1.58} 70B model supports a batch size that is up to 11 times larger than that of the \text{LLaMA LLM}, resulting in an 8.9-fold improvement in throughput.

\textbf{\text{BitNet b1.58} introduces a novel scaling law that relates model quality to inference efficiency.} Drawing on insights from Figure~\ref{fig:latency-memory-energy} and \ref{fig:energy}, a rough equivalence emerges between model sizes in the 1.58-bit and 16-bit settings:

\begin{itemize}
    \item The 13B BitNet b1.58 exhibits greater efficiency—regarding latency, memory footprint, and energy expenditure—than a 3B FP16 LLM.
    \item The 30B BitNet b1.58 surpasses a 7B FP16 LLM in latency, memory consumption, and energy use.
    \item The 70B BitNet b1.58 is more efficient than a 13B FP16 LLM, considering latency, memory utilization, and energy requirements.
\end{itemize}

\paragraph{Training with 2T Tokens}

For large language models, the amount of training tokens directly influences downstream performance. To evaluate the scaling behavior of \text{BitNet b1.58} with respect to training tokens, we pre-trained a \text{BitNet b1.58} variant using 2T tokens, following the same data protocol as StableLM-3B~\cite{StableLM-3B-4E1T}, which currently sets the open-source standard for 3B models. Both models underwent evaluation on a suite of benchmarks: Winogrande~\cite{winoGrande}, PIQA~\cite{piqa}, SciQ~\cite{sciq}, LAMBADA~\cite{lambada}, and ARC-easy~\cite{arc}. We present zero-shot results in Table~\ref{tab:2t}; when tasks report both accuracy and normalized accuracy, we average these metrics. StableLM 3B results for 2T tokens are referenced from its published report. Our analysis reveals that \text{BitNet b1.58} attains higher performance across all tasks, demonstrating that LLMs with 1.58-bit precision also possess strong generalization potential.

\section{Discussion and Future Work}

\textbf{1-bit Mixture-of-Experts (MoE) LLMs}

Mixture-of-Experts (MoE) models have established themselves as a resource-efficient methodology within large language models (LLMs). Although they offer a substantial reduction in computational FLOPs, MoE models are often constrained by significant memory demands and intensive inter-chip communication, which impede their practical adoption. These limitations can be mitigated by leveraging 1.58-bit LLMs. With a lower memory requirement, fewer hardware devices are needed for deployment. Additionally, the network traffic for transferring activations is significantly lessened. Ultimately, if the model fits onto a single chip, the communication overhead can be eliminated entirely.

\textbf{Native Support of Long Sequence in LLMs}

Supporting extended input sequences is now an essential requirement for LLMs. A central obstacle in accommodating long context lengths is the substantial memory required for KV cache storage during inference. \text{BitNet b1.58} offers notable improvements in this respect by cutting activation memory from 16 bits to 8 bits, effectively enabling the model to process twice as much context using the same memory resources. There remains the possibility to compress this even further to 4 bits or less, with minimal to no information loss for 1.58-bit LLMs. Exploration of such strategies will be part of our future investigations.

\textbf{LLMs on Edge and Mobile}

Deploying 1.58-bit LLMs can significantly enhance the functionality of language models on edge and mobile hardware, which typically suffer from limited memory and compute capabilities. By lowering memory and energy requirements, 1.58-bit LLMs make feasible a range of applications previously unattainable on such devices. This not only augments the practical utility of edge and mobile platforms but also unlocks new opportunities for LLM-driven technologies. Furthermore, the 1.58-bit architecture aligns particularly well with CPUs—prevalent in these environments—enabling \text{BitNet b1.58} to run efficiently, thereby further boosting system performance and expanding potential applications.

\textbf{New Hardware for 1-bit LLMs}

Developments like Groq\footnote{\url{https://groq.com/}} have illustrated the promise of developing specialized hardware (e.g., LPUs) tailored for LLM acceleration. Building upon this direction, there is a strong case and urgent need to pursue new hardware and systems specifically tailored for 1-bit LLMs, to take full advantage of the computational efficiencies made possible by BitNet~\cite{bitnet}.